\hypertarget{chapter-06-page-138}{}
\subsection{注记与进一步结果}
\label{sec:ch6-notes-further-results}

\subsubsection{参考文献}
\label{subsec:ch6-references}

关于 Hardy 空间的发展, 特别是 $H^1$ 空间, 见下文第~\ref{subsec:ch6-hp-spaces}~节. E. M. Stein 在
\MBUnresolvedReference{citation}{[Classes $H^p$, multiplicateurs et fonctions de Littlewood--Paley, C. R. Acad. Sci. Paris 263 (1966), 716--719, 780--781]} 中首先证明了奇异积分和乘子从 $H^1$ 到 $L^1$ 有界.

BMO 空间由 F. John 和 L. Nirenberg 在研究偏微分方程时引入; 见 \MBUnresolvedReference{citation}{[On functions of bounded mean oscillation, Comm. Pure Appl. Math. 14 (1961), 415--426]}. 他们在该文中证明了定理~\ref{thm:ch6-john-nirenberg}. 奇异积分把 $L^\infty$ 映入 BMO 这一事实最早由 S. Spanne, J. Peetre 和 E. M. Stein 观察到; 分别见 \MBUnresolvedReference{citation}{[Sur l'interpolation entre les espaces $L_k^{p,\Phi}$, Ann. Scuola Norm. Sup. Pisa 20 (1966), 625--648]}, \MBUnresolvedReference{citation}{[On convolution operators leaving $L^{p,\lambda}$ spaces invariant, Ann. Mat. Pura Appl. 72 (1966), 295--304]} 以及 \MBUnresolvedReference{citation}{[Singular integrals, harmonic functions, and differentiability properties of functions of several variables, Singular Integrals, Proc. Sympos. Pure Math. X, pp. 316--335, Amer. Math. Soc., Providence, 1967]}. 锐极大函数由 C. Fefferman 和 E. M. Stein 引入; 见 \MBUnresolvedReference{citation}{[$H^p$ spaces of several variables, Acta Math. 129 (1972), 137--193]}. 他们在同一篇论文中证明了定理~\ref{thm:ch6-linfty-bmo-interpolation}.

\subsubsection{$H^p$ 空间}
\label{subsec:ch6-hp-spaces}

Hardy 空间最初由 G. H. Hardy 作为复分析的一部分加以研究; 见 \MBUnresolvedReference{citation}{[The mean value of the modulus of an analytic function, Proc. London Math. Soc. 14 (1914), 269--277]}. 若单位圆盘 $\mathbb D=\{z\in\mathbb C:|z|<1\}$ 中的解析函数 $F$ 满足
\[
 \sup_{0<r<1}\int_{-\pi}^\pi|F(re^{i\theta})|^p\,d\theta<\infty,
 \qquad 0<p<\infty,
\]
则称 $F\in H^p(\mathbb D)$. 若上半平面 $\mathbb R_+^2=\{z=x+iy:y>0\}$ 中的解析函数 $F$ 满足
\[
 \sup_{y>0}\int_{\mathbb R}|F(x+iy)|^p\,dx<\infty,
\]
则称 $F\in H^p(\mathbb R_+^2)$.

当 $p>1$ 时, 可以证明 $H^p$ 等于如下解析函数类: 其实部是在圆周 $\mathbb T$ 或实直线 $\mathbb R$ 上的某个 $L^p$ 函数的 Poisson 积分. 因而可以分别把 $H^p(\mathbb D)$ 与 $L^p(\mathbb T)$, $H^p(\mathbb R_+^2)$ 与 $L^p(\mathbb R)$ 认同. 当 $p\le1$ 时这种认同不成立. $H^p(\mathbb D)$ 和 $H^p(\mathbb R_+^2)$ 中的函数具有以因子分解定理为基础的丰富结构理论, 这些定理允许把函数分解成容易理解的分量. 关于这种观点的系统论述, 见 Koosis \MBUnresolvedReference{citation}{[11]} 或 Rudin \MBUnresolvedReference{citation}{[14]}.

遗憾的是, 利用多复变函数理论不能把这些结果推广到高维. E. M. Stein 和 G. Weiss 在 \MBUnresolvedReference{citation}{[On the theory of harmonic functions of several variables, I: The theory of $H^p$ spaces, Acta Math. 103 (1960), 25--62]} 中通过考察满足广义 Cauchy--Riemann 方程的向量值函数定义了 $H^p(\mathbb R_+^{n+1})$; 另见 Stein \MBUnresolvedReference{citation}{[15, Chapter 7]}. 该推广只在 $p>1-1/n$ 时有效, 因而包括 $H^1$. 通过取边界值, 这个 $H^1$ 空间同构于 $L^1(\mathbb R^n)$ 中所有 Riesz 变换也属于 $L^1(\mathbb R^n)$ 的函数空间, 即第~\ref{sec:ch6-atomic-h1}~节定义的 $H^1(\mathbb R^n)$.

\hypertarget{chapter-06-page-139}{}
D. Burkholder, R. Gundy 和 M. Silverstein 在 \MBUnresolvedReference{citation}{[A maximal characterization of the class $H^p$, Trans. Amer. Math. Soc. 157 (1971), 137--153]} 中证明了建立纯实变量 $H^p$ 理论的一个关键结果. 设 $F=u+iv$ 是上半平面中的解析函数. 他们证明 $F\in H^p$ 当且仅当其实部的非切向极大函数
\[
 u_a^*(x_0)=\sup\{|u(x,y)|:(x,y)\in\Gamma_a(x_0)\},
\]
属于 $L^p(\mathbb R)$, 其中
\[
 \Gamma_a(x_0)=\{(x,y)\in\mathbb R_+^2:|x-x_0|<ay\}.
\]
由于 $u(x,y)=P_y*f(x)$ 是其边界值 $f$ 的 Poisson 积分, 这就刻画了 $H^p(\mathbb R_+^2)$ 中函数实部在 $\mathbb R$ 上的限制; 非切向逼近见第~\ref{subsec:ch2-nontangential-poisson-maximal}~节. 这个实直线上的函数空间通常记作 $H^p(\mathbb R)$, 有时也记作 $\operatorname{Re}H^p(\mathbb R)$. 当 $p=1$ 时, 它等价于第~\ref{sec:ch6-atomic-h1}~节定义的 $H^1(\mathbb R)$. 原证明使用 Brownian 运动; P. Koosis 给出了复分析证明, 见 \MBUnresolvedReference{citation}{[Sommabilité de la fonction maximale et appartenance à $H_1$, C. R. Acad. Sci. Paris 28 (1978), 1041--1043]} 及 \MBUnresolvedReference{citation}{[11, Chapter 8]}.

该结果的关键在于 $u$ 是调和函数这一事实并不重要: 若把 Poisson 核换成任意由积分非零的 $\phi\in\mathcal S$ 构造的近似恒等核, 边界值仍具有同样的刻画. 这给出高维 $H^p$ 的一个简洁定义. 给定积分非零的 $\phi\in\mathcal S(\mathbb R^n)$, 若
\[
 M_\phi f(x_0)=\sup\{|\phi_y*f(x)|:(x,y)\in\Gamma_a(x_0)\}
\]
属于 $L^p(\mathbb R^n)$, 就称 $f\in H^p(\mathbb R^n)$. 所得空间与 $\phi$ 无关. 实际上, 可以用所谓的 grand 极大函数
\[
 \mathcal Mf(x)=\sup_\phi M_\phi f(x)
\]
代替 $M_\phi$, 其中上确界取遍所有这样的 $\phi$. 当 $p>1-1/n$ 时, 通过取边界值, 该空间同构于上面 Stein 和 Weiss 的 $H^p$ 空间. 这种方法见上述 Fefferman 和 Stein 的论文, 另见 Stein \MBUnresolvedReference{citation}{[17, Chapter 3]}.

\hypertarget{chapter-06-page-140}{}
$H^p(\mathbb R)$ 中函数的原子分解归功于 R. Coifman, 见 \MBUnresolvedReference{citation}{[A real variable characterization of $H^p$, Studia Math. 51 (1974), 269--274]}; R. Latter 将其推广到高维, 见 \MBUnresolvedReference{citation}{[A decomposition of $H^p(\mathbb R^n)$ in terms of atoms, Studia Math. 62 (1978), 93--101]}. 另见 R. Latter 和 A. Uchiyama 的 \MBUnresolvedReference{citation}{[The atomic decomposition for parabolic $H^p$ spaces, Trans. Amer. Math. Soc. 253 (1979), 391--398]}. 为简单起见, 我们以原子分解作为定义 $H^1(\mathbb R^n)$ 的基础, 即使用称为 $H^1_{\mathrm{at}}(\mathbb R^n)$ 的空间. R. Coifman 和 G. Weiss 用这些方法在齐型空间上定义了 $H^p$ 空间, 见 \MBUnresolvedReference{citation}{[Extensions of Hardy spaces and their use in Analysis, Bull. Amer. Math. Soc. 83 (1977), 569--645]} 以及第~\ref{subsec:ch5-spaces-homogeneous-type}~节. 他们的定义只对 $p_0<p\le1$ 成立, 其中 $p_0$ 是依赖于所考察空间的常数.

Shanzhen Lu 的专著 \MBUnresolvedReference{citation}{[Four Lectures on Real $H^p$ Spaces, World Scientific, Singapore, 1995]} 论述了 $H^p$ 空间的实变量理论及其在 Fourier 分析和逼近论中的应用.

最后说明用 Lusin 面积积分刻画 $H^p$ 空间的另一种方法. 若 $u$ 是 $\mathbb R_+^{n+1}$ 上的调和函数, 令
\[
 S(u)(x_0)=\left(\int_{\Gamma_a^h(x_0)}|\nabla u|^2y^{1-n}\,dy\,dx\right)^{1/2},
\]
其中
\[
 \Gamma_a^h(x_0)=\{(x,y)\in\mathbb R_+^{n+1}:|x-x_0|<ay,\ 0<y<h\},
\]
且
\[
 |\nabla u|^2=\left|\frac{\partial u}{\partial y}\right|^2
 +\sum_{j=1}^n\left|\frac{\partial u}{\partial x_j}\right|^2.
\]
之所以称为面积积分, 是因为当 $n=1$ 时, $S(u)^2$ 等于实部为 $u$ 的解析函数 $F$ 将 $\Gamma_a^h(x_0)$ 映出的区域面积, 并计入重数. 此时 $u\in H^p(\mathbb R_+^{n+1})$ 当且仅当 $S(u)\in L^p(\mathbb R^n)$ 且当 $y\to\infty$ 时 $u(x,y)\to0$. 也可以用其径向类似物, 一个 Littlewood--Paley 型函数
\[
 g(u)(x)=\left(\int_0^\infty|\nabla u(x,y)|^2y\,dy\right)^{1/2},
\]
代替面积积分. 这一刻画归功于 A. P. Calderón, 见 \MBUnresolvedReference{citation}{[Commutators of singular integral operators, Proc. Nat. Acad. Sci. U.S.A. 53 (1965), 1092--1099]}. 另见 García-Cuerva 和 Rubio de Francia \MBUnresolvedReference{citation}{[6]}, Stein \MBUnresolvedReference{citation}{[15]}, Zygmund \MBUnresolvedReference{citation}{[21]} 以及上述 Fefferman 和 Stein 的论文.

\hypertarget{chapter-06-page-141}{}
\subsubsection{对偶性}
\label{subsec:ch6-duality}

从 $H^1$ 到 $L^1$ 有界和从 $L^\infty$ 到 BMO 有界这两个结果彼此对偶. 这是 C. Fefferman 的一个深刻结果的推论; 该结果公布于 \MBUnresolvedReference{citation}{[Characterizations of bounded mean oscillation, Bull. Amer. Math. Soc. 77 (1971), 587--588]}.

\begin{Theorem}
\label{thm:ch6-h1-dual-bmo}
$H^1(\mathbb R^n)$ 的对偶空间是 BMO.
\end{Theorem}

上述 Fefferman 和 Stein 的论文给出该结果的两个证明, García-Cuerva 和 Rubio de Francia \MBUnresolvedReference{citation}{[6]} 中给出三个证明. 当 $p<1$ 时, $H^p$ 空间的对偶是 Lipschitz 空间. 在单位圆周上, 该结果归功于 P. Duren, B. Romberg 和 A. Shields, 见 \MBUnresolvedReference{citation}{[Linear functionals on $H^p$ spaces with $0<p<1$, J. Reine Angew. Math. 238 (1969), 32--60]}; 在 $\mathbb R^n$ 上归功于 T. Walsh, 见 \MBUnresolvedReference{citation}{[The dual of $H^p(\mathbb R_+^{n+1})$ for $p<1$, Can. J. Math. 25 (1973), 567--577]}. 另见 \MBUnresolvedReference{citation}{[6, Chapter 3]}.

\subsubsection{BMO 上的 Hardy--Littlewood 极大函数}
\label{subsec:ch6-maximal-on-bmo}

极大函数在 BMO 上有界: 若 $f\in\mathrm{BMO}$, 则 $Mf\in\mathrm{BMO}$. 该结果首先由 C. Bennett, R. A. DeVore 和 R. Sharpley 证明, 见 \MBUnresolvedReference{citation}{[Weak $L^\infty$ and BMO, Ann. of Math. 113 (1981), 601--611]}. F. Chiarenza 和 M. Frasca, 以及 D. Cruz-Uribe 和 C. J. Neugebauer 分别给出较简单的证明; 见 \MBUnresolvedReference{citation}{[Morrey spaces and Hardy--Littlewood maximal function, Rend. Mat. Appl., Series 7, 7 (1981), 273--279]} 和 \MBUnresolvedReference{citation}{[The structure of the reverse Hölder classes, Trans. Amer. Math. Soc. 347 (1995), 2941--2960]}. 另见 Torchinsky \MBUnresolvedReference{citation}{[19, p. 204]}.

\subsubsection{另一个插值结果}
\label{subsec:ch6-another-interpolation}

应用 Marcinkiewicz 插值定理时, 常常使用算子在 $L^\infty$ 上有界或满足弱型 $(1,1)$ 不等式的事实. 定理~\ref{thm:ch6-linfty-bmo-interpolation} 表明, 可以用算子把 $L^\infty$ 映入 BMO 来代替第一个假设. 也可以用 $H^1$ 代替弱型 $(1,1)$ 不等式.

\begin{Theorem}
\label{thm:ch6-h1-weak-interpolation}
设 $T$ 是次线性算子, 它把 $H^1$ 映入 $L^1$, 并且对某个 $1<p_1\le\infty$, $T$ 是弱型 $(p_1,p_1)$ 的. 则对每个 $1<p<p_1$, $T$ 在 $L^p$ 上有界.
\end{Theorem}

该结果可以推广到其他 $H^p$ 空间. 细节与参考文献见 García-Cuerva 和 Rubio de Francia \MBUnresolvedReference{citation}{[6, Chapter 3]}.

\subsubsection{分数积分与 Sobolev 空间}
\label{subsec:ch6-fractional-sobolev}

对分数积分算子 $I_\alpha$, 临界指标是 $p=n/\alpha$. 我们本来会期望 $I_\alpha$ 把 $L^{n/\alpha}$ 映入 $L^\infty$, 但事实并非如此. 取而代之的是定理~\ref{thm:ch6-linfty-to-bmo} 的类似结果.

\hypertarget{chapter-06-page-142}{}
\begin{Theorem}
\label{thm:ch6-fractional-integral-to-bmo}
若 $f\in L^{n/\alpha}$, 则 $I_\alpha f\in\mathrm{BMO}$.
\end{Theorem}

该定理由联系分数积分, 锐极大函数和分数极大函数的一个不等式立即推出.

\begin{Theorem}
\label{thm:ch6-fractional-sharp-maximal-comparison}
存在正常数 $C_1,C_2$, 使每个局部可积函数 $f$ 都满足
\[
 C_1M_\alpha f(x)\le M^\#(I_\alpha f)(x)\le C_2M_\alpha f(x).
\]
\end{Theorem}

这两个结果都归功于 D. R. Adams, 见 \MBUnresolvedReference{citation}{[A note on Riesz potentials, Duke Math. J. 42 (1975), 765--778]}.

当 $p=1$ 时, 除了定理~\ref{thm:ch4-hardy-littlewood-sobolev} 中的弱型 $(1,n/(n-\alpha))$ 不等式, 还有推论~\ref{cor:ch6-singular-integral-h1-to-l1} 的类似结果:
\[
 \|I_\alpha f\|_{n/(n-\alpha)}\le C\|f\|_{H^1}.
\]
该结果归功于 Stein 和 Weiss, 见上文所引论文. 对所有 $H^p$ 空间都有类似不等式, 见 S. Krantz 的 \MBUnresolvedReference{citation}{[Fractional integration on Hardy Spaces, Studia Math. 73 (1982), 87--94]}.

对 Sobolev 空间, 相应的临界定理应为: 若 $f\in L_k^{n/k}$, 则 $f\in\mathrm{BMO}$. 但这只在 $k=1$ 时已知, 见 A. Cianchi 和 L. Pick 的 \MBUnresolvedReference{citation}{[Sobolev embeddings into BMO, VMO and $L^\infty$, Ark. Mat. 36 (1998), 317--340]}. 一般情形可用一个指数可积性结果代替, 比较推论~\ref{cor:ch6-bmo-exponential-integrability}.

\begin{Theorem}
\label{thm:ch6-trudinger-exponential-integrability}
给定 $0<\alpha<n$ 和任意立方体 $Q$, 存在常数 $C_1,C_2$, 使当 $f\in L^{n/\alpha}(Q)$ 时,
\[
 \frac1{|Q|}\int_Q
 \exp\left(\left|\frac{I_\alpha f}{C_1\|f\|_{L^{n/\alpha}(Q)}}\right|^{p'}\right)
 \le C_2.
\]
\end{Theorem}

当 $\alpha=1$ 时, 定理~\ref{thm:ch6-trudinger-exponential-integrability} 由 N. S. Trudinger 证明, 见 \MBUnresolvedReference{citation}{[On imbeddings into Orlicz spaces and some applications, J. Math. Mech. 17 (1967), 473--483]}; 一般情形由 R. S. Strichartz 证明, 见 \MBUnresolvedReference{citation}{[A note on Trudinger's extension of Sobolev's inequalities, Indiana Univ. Math. J. 21 (1972), 841--842]}. J. Moser 和 D. R. Adams 得到了最佳常数, 分别见 \MBUnresolvedReference{citation}{[A sharp form of an inequality by N. Trudinger, Indiana Univ. Math. J. 20 (1971), 1077--1092]} 与 \MBUnresolvedReference{citation}{[A sharp inequality of J. Moser for higher order derivatives, Ann. of Math. 128 (1988), 385--398]}. 另见 W. Ziemer 的 \MBUnresolvedReference{citation}{[Weakly Differentiable Functions, Springer-Verlag, New York, 1989]} 以及 D. R. Adams 和 L. Hedberg 的 \MBUnresolvedReference{citation}{[Function Spaces and Potential Theory, Springer-Verlag, Berlin, 1996]}.

\hypertarget{chapter-06-page-143}{}
\subsubsection{消失平均振荡}
\label{subsec:ch6-vmo}

给定 $f\in\mathrm{BMO}$, 若
\[
 \frac1{|Q|}\int_Q|f(x)-f_Q|\,dx
\]
当 $|Q|$ 趋于 $0$ 或 $\infty$ 时都趋于零, 就说 $f$ 具有消失平均振荡, 记作 $f\in\mathrm{VMO}$. 与 BMO 一样, VMO 中也有无界函数. 例如
\[
 f(x)=
 \begin{cases}
  \log\log(1/|x|),&|x|\le1/e,\\
  0,&|x|>1/e,
 \end{cases}
\]
属于 VMO. 事实上, 可以证明在某种意义下它具有 VMO 函数所能达到的最大增长速度. 显然, 若 $f\in C_0$, 即 $f$ 连续且在无穷远处趋于零, 则 $f\in\mathrm{VMO}$. 进一步, VMO 是 $C_0$ 在 BMO 中的闭包.

VMO 可把定理~\ref{thm:ch6-linfty-to-bmo} 加强如下: 给定奇异积分算子 $T$, 若 $f\in C_0$, 则 $Tf\in\mathrm{VMO}$.

VMO 空间由 D. Sarason 引入, 见 \MBUnresolvedReference{citation}{[Functions of vanishing mean oscillation, Trans. Amer. Math. Soc. 207 (1975), 391--405]}. 另见 R. Coifman 和 G. Weiss 的综述 \MBUnresolvedReference{citation}{[Extensions of Hardy spaces and their uses in analysis, Bull. Amer. Math. Soc. 83 (1977), 569--645]}.

\subsubsection{交换子与 BMO}
\label{subsec:ch6-commutators-bmo}

可以用奇异积分给出 BMO 的另一个刻画. 设 $T$ 是奇异积分, 对局部可积函数 $b$, 令 $M_b$ 是由 $b$ 逐点乘法给出的算子: $M_bf(x)=b(x)f(x)$. 定义交换子 $[b,T]=M_bT-TM_b$. 若 $Tf=K*f$, 则对 $f\in C_c^\infty$,
\[
 [b,T]f(x)=\int_{\mathbb R^n}(b(x)-b(y))K(x-y)f(y)\,dy,
 \qquad x\notin\operatorname{supp}(f).
\]
显然, 若 $b\in L^\infty$, 则 $[b,T]$ 在 $L^p$, $1<p<\infty$, 上有界. 因为 $M_bT$ 与 $TM_b$ 都有界恰好发生在 $b$ 有界时, 猜测 $[b,T]$ 也只有在此时才有界是合理的. 但是, 两项之间的抵消使更弱的条件已经充分.

\begin{Theorem}
\label{thm:ch6-commutator-bmo-characterization}
设奇异积分 $T$ 的核是标准核. 则交换子 $[b,T]$ 在 $L^p$, $1<p<\infty$, 上有界, 当且仅当 $b\in\mathrm{BMO}$.
\end{Theorem}

交换子由 R. Coifman, R. Rochberg 和 G. Weiss 引入; 见 \MBUnresolvedReference{citation}{[Factorization theorems for Hardy spaces in several variables, Ann. of Math. 103 (1976), 611--635]}. 他们用交换子把经典 $H^p$ 理论推广到高维, 证明了 $b\in\mathrm{BMO}$ 对 $[b,T]$ 有界是充分的, 并证明了部分逆命题.

\hypertarget{chapter-06-page-144}{}
完整逆命题归功于 S. Janson, 见 \MBUnresolvedReference{citation}{[Mean oscillation and commutators of singular integral operators, Ark. Mat. 16 (1978), 263--270]}. 该文还包含 J. O. Strömberg 给出的一个更简单的充分性证明, Torchinsky \MBUnresolvedReference{citation}{[19, p. 418]} 重述了该证明.

交换子 $[b,T]$ 比相应的奇异积分更奇异. 特别地, 它不满足相应的弱型 $(1,1)$ 估计. 但下述较弱结果成立.

\begin{Theorem}
\label{thm:ch6-commutator-endpoint}
给定奇异积分 $T$ 和 $b\in\mathrm{BMO}$, 则
\[
 |\{x\in\mathbb R^n:|[b,T]f(x)|>\lambda\}|
 \le C\|b\|_*\int_{\mathbb R^n}\frac{|f(y)|}{\lambda}
 \left(1+\log^+\left(\frac{|f(y)|}{\lambda}\right)\right)\,dy.
\]
\end{Theorem}

定理~\ref{thm:ch6-commutator-endpoint} 归功于 C. Pérez, 见 \MBUnresolvedReference{citation}{[Endpoint estimates for commutators of singular integral operators, J. Funct. Anal. 128 (1995), 163--185]}. 他还在该文中讨论了 $[b,T]$ 在 $H^1$ 上的有界性.

A. Uchiyama 在 \MBUnresolvedReference{citation}{[On the compactness of operators of Hankel type, Tôhoku Math. J. 30 (1978), 163--171]} 中证明了交换子 $[b,T]$ 是紧算子当且仅当 $b\in\mathrm{VMO}$.
